\documentclass[a4paper,12pt]{article}
\usepackage{listings}
\usepackage{xcolor}
\usepackage{geometry}
\usepackage{hyperref}
\geometry{margin=1in}
\hypersetup{
    colorlinks=true,
    linkcolor=blue,
    urlcolor=blue,
}

\definecolor{codegray}{rgb}{0.95,0.95,0.95}
\lstdefinestyle{mystyle}{
    backgroundcolor=\color{codegray},
    basicstyle=\ttfamily\footnotesize,
    breaklines=true,
    frame=single,
    keywordstyle=\color{blue},
    commentstyle=\color{green!50!black},
    stringstyle=\color{red!60!black},
    showstringspaces=false
}
\lstset{style=mystyle}

\title{SYCL Foundation & CUTLASS-SYCL Roadmap}
\author{Rishi Yadav}
\date{\today}

\begin{document}

\maketitle

\tableofcontents
\newpage

\section{Introduction}
This document provides a comprehensive roadmap to learn SYCL and transition to CUTLASS-SYCL. It covers the core SYCL concepts, essential lessons, and their relevance for building high-performance GEMM kernels using CUTLASS in SYCL.

\section{SYCL Foundation}

\subsection{Lesson 1: Hello SYCL}
\textbf{Concepts:} Queue, device selection, basic host code.

\begin{lstlisting}[language=C++]
#include <sycl/sycl.hpp>
using namespace sycl;

int main() {
    queue q;
    std::cout << "Hello, SYCL from: "
              << q.get_device().get_info<info::device::name>()
              << std::endl;
    return 0;
}
\end{lstlisting}

\textbf{Key Points:}
\begin{itemize}
    \item \texttt{queue} selects the device.
    \item Host code can query device info.
\end{itemize}

\subsection{Lesson 2: First Kernel}
\textbf{Concepts:} Parallel work-items, device code.

\begin{lstlisting}[language=C++]
#include <sycl/sycl.hpp>
#include <vector>
#include <iostream>
using namespace sycl;

int main() {
    queue q;
    const int N = 16;
    std::vector<int> data(N);

    {
        buffer<int,1> buf(data.data(), range<1>(N));
        q.submit([&](handler& h) {
            auto acc = buf.get_access<access::mode::write>(h);
            h.parallel_for(range<1>(N), [=](id<1> i){
                acc[i] = i[0];
            });
        }).wait();
    }

    for(int i=0;i<N;i++)
        std::cout << "Hello from work-item " << data[i] << "\n";

    return 0;
}
\end{lstlisting}

\subsection{Lesson 3: Buffers \& Accessors}
\textbf{Concepts:} Host ↔ Device memory, read/write access.

\begin{lstlisting}[language=C++]
#include <sycl/sycl.hpp>
#include <vector>
#include <iostream>
using namespace sycl;

int main() {
    queue q;
    const int N = 16;
    std::vector<int> A(N), B(N), C(N);

    for(int i=0;i<N;i++){
        A[i]=i; B[i]=2*i;
    }

    {
        buffer<int,1> bufA(A.data(), range<1>(N));
        buffer<int,1> bufB(B.data(), range<1>(N));
        buffer<int,1> bufC(C.data(), range<1>(N));

        q.submit([&](handler& h){
            auto accA = bufA.get_access<access::mode::read>(h);
            auto accB = bufB.get_access<access::mode::read>(h);
            auto accC = bufC.get_access<access::mode::write>(h);

            h.parallel_for(range<1>(N), [=](id<1> i){
                accC[i] = accA[i] + accB[i];
            });
        }).wait();
    }

    for(int i=0;i<N;i++)
        std::cout << A[i] << " + " << B[i] << " = " << C[i] << "\n";

    return 0;
}
\end{lstlisting}

\subsection{Lesson 4: ND-Range \& Work-Groups}
\textbf{Concepts:} Global and local size, 2D indexing.

\begin{lstlisting}[language=C++]
#include <sycl/sycl.hpp>
#include <iostream>
#include <vector>
using namespace sycl;

int main() {
    queue q;
    const int N = 4;
    std::vector<int> A(N*N), B(N*N), C(N*N);

    for(int i=0;i<N*N;i++){
        A[i]=i; B[i]=i*2;
    }

    {
        buffer<int,2> bufA(A.data(), range<2>(N,N));
        buffer<int,2> bufB(B.data(), range<2>(N,N));
        buffer<int,2> bufC(C.data(), range<2>(N,N));

        q.submit([&](handler& h){
            auto accA = bufA.get_access<access::mode::read>(h);
            auto accB = bufB.get_access<access::mode::read>(h);
            auto accC = bufC.get_access<access::mode::write>(h);

            h.parallel_for(nd_range<2>(range<2>(N,N), range<2>(2,2)),
                [=](nd_item<2> item){
                    int row = item.get_global_id(0);
                    int col = item.get_global_id(1);
                    accC[row][col] = accA[row][col] + accB[row][col];
            });
        }).wait();
    }

    for(int i=0;i<N;i++){
        for(int j=0;j<N;j++)
            std::cout << C[i*N + j] << " ";
        std::cout << "\n";
    }

    return 0;
}
\end{lstlisting}

\subsection{Lesson 5: Local Memory \& Barriers}
\textbf{Concepts:} Fast shared memory, synchronization.

\begin{lstlisting}[language=C++]
// Local memory example included in previous lessons
\end{lstlisting}

\subsection{Lesson 6: Sub-Groups \& SIMD}
\textbf{Concepts:} Vectorized computation, warp-level operations.

\begin{lstlisting}[language=C++]
// Sub-group example included in previous lessons
\end{lstlisting}

\subsection{Lesson 7: Unified Shared Memory (USM)}
\textbf{Concepts:} Pointer-style memory, host/device shared pointers.

\begin{lstlisting}[language=C++]
#include <sycl/sycl.hpp>
using namespace sycl;

int main() {
    queue q;
    const int N = 16;
    int *A = malloc_shared<int>(N, q);
    int *B = malloc_shared<int>(N, q);
    int *C = malloc_shared<int>(N, q);

    for(int i=0;i<N;i++){
        A[i] = i;
        B[i] = i*2;
    }

    q.parallel_for(range<1>(N), [=](id<1> i){
        C[i] = A[i] + B[i];
    }).wait();

    for(int i=0;i<N;i++)
        std::cout << C[i] << " ";
    std::cout << "\n";

    free(A,q); free(B,q); free(C,q);
}
\end{lstlisting}

\section{CUTLASS-SYCL Roadmap}
\begin{itemize}
    \item Clone the repository: \texttt{git clone https://github.com/intel/cutlass-sycl.git}
    \item Build small GEMM kernels first (CPU → GPU/NPU)
    \item Map SYCL lessons to GEMM:
    \begin{itemize}
        \item Buffers/USM → matrices
        \item ND-Range → tiles/work-groups
        \item Local memory → shared tile storage
        \item Sub-groups → SIMD MAC operations
    \end{itemize}
    \item Experiment with tile sizes, local ranges, and vectorization
\end{itemize}

\section{Conclusion}
After completing SYCL lessons 1–7 and understanding ND-range, local memory, sub-groups, and USM, you are ready to dive into CUTLASS-SYCL. This foundation ensures you understand GEMM kernels, memory hierarchy, and device parallelism.

\end{document}
